\section{Descripción de la industria}

\subsection{Institucionalidad del sector}

\paragraph{La cadena productiva} La industria de los combustibles líquidos tiene cuatro etapas de producción: extracción, refinación, distribución mayorista y distribución minorista \parencite{FNE2025informe}. Siendo la distribución minorista la mitad de las ventas antuales de combustible del país. La ENAP es la única empresa que produce y refina petróleo crudo en Chile, insumo que es esencialmente importado siendo la producción nacional alrededor del $1\%$ con respecto a las importaciones. Si bien las grandes compañías mayoristas pueden abastecerse tanto de las refinerías de ENAP como importar directamente, las empresas de menor tamaño se abastecen únicamente de ENAP \parencite{FNE2025informe}.\footnote{El $12\%$ de la oferta nacional de gasolinas y el $73\%$ del diésel son importados} La consecuencia relevante para este trabajo es que el costo de adquisición de las EDS está anclado, para la generalidad del mercado, a una referencia única.

En términos de producto, las gasolinas de 93, 95 y 97 octanos y el petróleo diésel constituyen el objeto de análisis. El diésel explica algo más de la mitad de las ventas nacionales de combustibles líquidos, aunque con usos diversos en la industria, el comercio y el transporte, las gasolinas, que son consumidas casi exclusivamente por el transporte terrestre, representan en conjunto alrededor de un tercio, con la de 93 octanos como la de mayor peso \parencite{FNE2025informe}.

\paragraph{El costo mayorista y el MEPCO} El precio mayorista que enfrentan los distribuidores se construye sobre el precio de paridad de importación que fija semanalmente el Ministerio de Energía, previo informe de la Comisión Nacional de Energía (CNE), y sobre él operan los impuestos específicos de la ley N°18.502.

La ley N°20.765 de 2014 creó el Mecanismo de Estabilización de Precios de los Combustibles (MEPCO), que atenúa la transmisión de las variaciones internacionales al precio interno \parencite{LEY20765}.\footnote{El mecanismo opera a través de un componente variable que se suma o se resta al componente base del impuesto específico. Su regla de cálculo es explícita: se proyecta un ``precio base'' como, el precio que ENAP informará la semana siguiente suponiendo un componente variable nulo, incluidos IVA e impuesto específico base, y se lo compara con el precio efectivamente informado la semana anterior. Cuando la diferencia excede el $0{,}8\%$ del promedio de las dos últimas semanas del precio base de la gasolina de 93 octanos, el componente variable se ajusta para absorber el exceso.}  En la práctica, el mecanismo acota la variación semanal del precio mayorista en torno a un umbral con el objetivo de estabilizar el precio final.

Tres rasgos del MEPCO son importantes para la estrategia empírica. Primero, es una referencia nacional, el precio de paridad se fija una vez por semana y entra en vigencia el jueves siguiente, sin diferenciación geográfica. Segundo, es una referencia fijada por regla y publicada, no negociada bilateralmente, por lo que no puede responder a condiciones de competencia locales. Tercero, la propia ley precisa que los precios de referencia y de paridad ``serán mera referencia y no constituirán precios mínimos ni máximos de venta'' \parencite{LEY20765}. Es decir, el MEPCO regula el costo mayorista, no el precio minorista, que las EDS fijan libremente. Es esta combinación la que permite aproximar el margen de cada estación e identificar en él el componente competitivo del precio.

\paragraph{Estructura del mercado minorista} La distribución minorista está dominada por tres compañías integradas verticalmente. Al año 2023, Copec concentraba el $39\%$ de las EDS del país, Petrobras el $21\%$ y Shell el $16\%$, el $24\%$ restante corresponde a distribuidoras de bandera blanca, esto es, estaciones que operan bajo marca propia o sin marca, con redes reducidas, menores costos operativos y menores estándares de servicio \parencite{FNE2025informe}.

La integración vertical no siempre implica propiedad. Las mayoristas participan en el mercado minorista mediante redes de EDS propias o de terceros, y la relación contractual admite dos formas que difieren en un punto decisivo para este análisis. Bajo el modelo de consignación, el combustible sigue siendo propiedad de la mayorista, que determina el precio de venta al público, y la administradora de la EDS recibe como remuneración un margen por litro vendido. Bajo el modelo de concesión, en cambio, la administradora compra previamente el combustible y es ella quien fija el precio al público \parencite{FNE2025informe}. La coexistencia de ambos regímenes implica que una fracción no menor de las decisiones de precio observadas en los datos se toma a nivel de cadena y no de estación, lo que motiva incorporar efectos fijos de distribuidor interactuados con el tiempo en las especificaciones empíricas.

\paragraph{Homogeneidad del producto y transparencia de precios} Dos rasgos regulatorios definen el tipo de competencia posible en este mercado. El primero es la homogeneidad, los combustibles deben cumplir especificaciones de calidad fijadas por el Ministerio de Energía, de modo que dos productos del mismo octanaje son sustitutos perfectos con independencia de la marca que los comercialice \parencite{FNE2025informe}. El segundo es la transparencia. Desde 2012 la CNE administra la plataforma Bencina en Línea, y las EDS están obligadas
por norma a informar sus precios de venta cada vez que éstos varían, obligación sujeta a fiscalización y multas.

\paragraph{El carácter local de la competencia} La FNE ha definido de manera consistente el mercado relevante de producto como la distribución minorista de combustibles líquidos a través de EDS, segmentable según tipo de combustible. Respecto del mercado geográfico, ha sostenido que en el caso de las estaciones de servicio éste es \textcquote{FNE2021informe}{%
eminentemente local, con una ampliación del área geográfica de competencia a través de cadenas de sustitución, cuyo efecto se disipa en la medida que la distancia es mayor, por lo que existiría un punto en que la presión competitiva de una estación de servicio determinada sobre la otra deja de ser relevante}.
En la práctica ha considerado radios de tres a cinco kilómetros, sin descartar radios menores.

\paragraph{Barreras a la entrada} Tanto la FNE como el Tribunal de Defensa de la Libre Competencia han constatado en repetidas ocasiones la existencia de barreras a la entrada en este mercado. Éstas se refieren principalmente a la escasez de terrenos aptos para la instalación o arriendo de estaciones de servicio, dadas las exigencias de los planos reguladores en las zonas urbanas más saturadas, que son al mismo tiempo las de mayor atractivo comercial \parencite{FNE2021informe}. A ello se suman la regulación sectorial a lo largo de toda la cadena productiva, que involucra a distintas autoridades, y la inversión en infraestructura que requiere un potencial entrante \parencite{FNE2025informe}.

\subsection{Modelo de competencia espacial}

El marco natural para este mercado es el modelo de ciudad circular de \textcite{SALOP1979}.\footnote{El modelo de \textcite{SALOP1979} es una variante del de \textcite{HOTELLING1929}, del que hereda todo lo que esta aplicación necesita: consumidores distribuidos uniformemente sobre un espacio, demanda unitaria e inelástica, un costo de desplazamiento $c$ por unidad de distancia y ninguna preferencia entre vendedores que no sea el precio más el costo de llegar hasta ellos. Se prefiere la versión circular porque la recta de \textcite{HOTELLING1929} tiene extremos y, con ellos, posiciones asimétricas que impiden un equilibrio simétrico con más de dos firmas. La circunferencia elimina ese problema y permite obtener el resultado en el número de competidoras, $p = m + c/n$, que es justamente la estática comparativa que este trabajo contrasta.} Los combustibles son un bien homogéneo, los precios son observables y la regulación impide diferenciación relevante en el producto, de modo que la única fuente sustantiva de diferenciación es la ubicación. Un consumidor prefiere la EDS más cercana salvo que otra compense la distancia con un precio menor. Esta subsección recoge únicamente lo indispensable para plantear las hipótesis. Los resultados no son una adaptación del modelo sino el modelo mismo, y el Anexo~\ref{anexo:modelo} los desarrolla, indica la ecuación y la página de \textcite{SALOP1979} de las que proviene cada uno y discute las limitaciones del marco.

El mercado local de una estación se representa por una circunferencia de perímetro unitario sobre la que se distribuyen uniformemente $L$ consumidores de demanda unitaria y operan $n$ EDS equiespaciadas, de modo que la distancia entre dos contiguas es $1/n$. Un consumidor que se desplaza una distancia $x$ incurre en un costo $cx$, con $c>0$. Cada estación enfrenta un costo fijo $F$ y un costo marginal $m$ que, por ser el MEPCO una referencia nacional, es común a todas ellas en una fecha dada.\footnote{Si bien los costos mayoristas sí difieren en la realidad entre estaciones, se hace este supuesto de igual costo mayorista tanto para el modelo teórico como para la parte empírica.} El margen bruto
\begin{equation}
    \mu_i \equiv p_i - m
    \label{eq:def-margen}
\end{equation}
recoge entonces toda la variación de precios entre estaciones en un mismo instante y es, entendido así, el componente competitivo del precio. Como el número de competidoras dentro del radio en que la sustitución es efectiva difiere entre estaciones de una misma comuna, cada EDS enfrenta su propio $n_i$, y es esa variación la que el modelo traduce en variación de márgenes.

\begin{proposicion}[Margen de equilibrio]
\label{prop:margen}
En el equilibrio de Nash en precios de la región competitiva, con $n$ estaciones
equiespaciadas,
\begin{equation}
    \mu = \frac{c}{n} .
    \label{eq:margen}
\end{equation}
\end{proposicion}

\begin{corolario}[Efecto de una entrada]
\label{cor:entrada}
La llegada de una competidora al mercado local de la incumbente altera su margen
de equilibrio en
\begin{equation}
    \Delta\mu = \frac{c}{n+1} - \frac{c}{n} = -\,\frac{c}{n(n+1)} \;<\; 0 ,
    \label{eq:efecto}
\end{equation}
y, dado que $m$ no se altera, el precio cae en la misma magnitud.
\end{corolario}

\paragraph{Hipótesis contrastables} De \eqref{eq:margen} y \eqref{eq:efecto} se desprenden las tres hipótesis que la estrategia empírica somete a prueba.

\begin{hipotesis}[Efecto disciplinador]
\label{hip:nivel}
La entrada de un competidor reduce el precio y el margen de las incumbentes del
mercado local.
\end{hipotesis}
El signo de \eqref{eq:efecto} es inequívoco, y el efecto es permanente porque opera sobre el equilibrio y no sobre una desviación transitoria.

\begin{hipotesis}[Atenuación espacial]
\label{hip:distancia}
El efecto decrece en magnitud con la distancia entre la entrante y la incumbente, y se anula a partir de un umbral.
\end{hipotesis}
Si la entrante queda fuera del radio en que la sustitución es efectiva, la incumbente sigue enfrentando el mismo $n_i$ y el efecto es exactamente nulo. La predicción a nivel de estación es entonces discontinua; el perfil suave que se estima por anillos de distancia resulta de promediar ese umbral sobre estaciones con radios de mercado distintos.

\begin{hipotesis}[Heterogeneidad por margen previo]
\label{hip:margen}
El efecto es mayor sobre las incumbentes cuyo margen previo a la entrada era más alto en relación con su mercado local.
\end{hipotesis}
Por \eqref{eq:margen} el margen es un estadístico suficiente del aislamiento espacial de la estación. Invirtiendo esa expresión como $n = c/\mu$ y sustituyéndola en \eqref{eq:efecto} se obtiene $|\Delta\mu| = \mu^{2}/(c+\mu)$, creciente en el margen previo tanto en niveles como en términos relativos.
